\section{Splitting finite function powers}

\import{set/function-split.tex}
\import{set/natural-shift.tex}

The first block retains its indices. Translation identifies the second block
with its own finite ordinal, so concatenation does not encode dimensions by
nested pairs. Both maps are defined explicitly, including empty blocks.

\begin{definition}\label{finite_function_split}
    $\finiteFunctionSplit{m}{n}{A}
        =\productMap{\funs{m}{A}}{\funs{(m+n)\setminus m}{A}}{
            \identity{\funs{m}{A}}}{
            \functionReindex{(m+n)\setminus m}{n}{A}{\finiteIndexShift{m}{n}}}
            \circ\functionSplit{m}{(m+n)\setminus m}{A}$.
\end{definition}

\begin{definition}\label{finite_function_join}
    $\finiteFunctionJoin{m}{n}{A}
        =\functionJoin{m}{(m+n)\setminus m}{A}\circ
            \productMap{\funs{m}{A}}{\funs{n}{A}}{
                \identity{\funs{m}{A}}}{
                \functionReindex{n}{(m+n)\setminus m}{A}{\converse{\finiteIndexShift{m}{n}}}}$.
\end{definition}

\begin{theorem}\label{finite_function_split_bijection}
    Let $m,n$ be natural numbers.
    Then $\finiteFunctionSplit{m}{n}{A}\in
        \bijections{\funs{m+n}{A}}{\funs{m}{A}\times\funs{n}{A}}$.
\end{theorem}
\begin{proof}
    Let $D=(m+n)\setminus m$ and $e=\finiteIndexShift{m}{n}$.
    $m\subseteq m+n$ by \obv{natural_addition_left_subseteq_sum}.
    Hence $m\union D=m+n$ by \auto{setminus_partition}.
    $m$ is disjoint from $D$ by \auto{disjoint,setminus}.
    Hence $\functionSplit{m}{D}{A}\in
        \bijections{\funs{m+n}{A}}{\funs{m}{A}\times\funs{D}{A}}$ by
        \auto{function_split_bijection}.
    $e\in\bijections{n}{D}$ by \auto{finite_index_shift_bijection}.
    Hence $\functionReindex{D}{n}{A}{e}\in\bijections{\funs{D}{A}}{\funs{n}{A}}$ by
        \obv{function_reindex_bijection}.
    $\identity{\funs{m}{A}}\in\bijections{\funs{m}{A}}{\funs{m}{A}}$ by
        \obv{id_elem_bijections}.
    Hence $\productMap{\funs{m}{A}}{\funs{D}{A}}{
        \identity{\funs{m}{A}}}{\functionReindex{D}{n}{A}{e}}\in
        \bijections{\funs{m}{A}\times\funs{D}{A}}{\funs{m}{A}\times\funs{n}{A}}$ by
        \obv{product_map_bijection}.
    Follows by \auto{finite_function_split,bijection_circ}.
\end{proof}

\begin{theorem}\label{finite_function_split_converse}
    Let $m,n$ be natural numbers.
    Then $\converse{\finiteFunctionSplit{m}{n}{A}}=\finiteFunctionJoin{m}{n}{A}$.
\end{theorem}
\begin{proof}
    Let $D=(m+n)\setminus m$ and $e=\finiteIndexShift{m}{n}$.
    $m$ is disjoint from $D$ by \auto{disjoint,setminus}.
    $\converse{\functionSplit{m}{D}{A}}=\functionJoin{m}{D}{A}$ by
        \obv{function_split_converse}.
    $e\in\bijections{n}{D}$ by \auto{finite_index_shift_bijection}.
    $\functionReindex{D}{n}{A}{e}\in\bijections{\funs{D}{A}}{\funs{n}{A}}$ by
        \obv{function_reindex_bijection}.
    $\converse{\functionReindex{D}{n}{A}{e}}=\functionReindex{n}{D}{A}{\converse{e}}$ by
        \obv{function_reindex_converse}.
    $\identity{\funs{m}{A}}\in\bijections{\funs{m}{A}}{\funs{m}{A}}$ by
        \obv{id_elem_bijections}.
    $\converse{\productMap{\funs{m}{A}}{\funs{D}{A}}{
        \identity{\funs{m}{A}}}{\functionReindex{D}{n}{A}{e}}}
        =\productMap{\funs{m}{A}}{\funs{n}{A}}{
            \identity{\funs{m}{A}}}{\functionReindex{n}{D}{A}{\converse{e}}}$ by
        \auto{product_map_converse,converse_identity}.
    Follows by \auto{finite_function_split,finite_function_join,circ_converse}.
\end{proof}

\begin{corollary}\label{finite_function_join_bijection}
    Let $m,n$ be natural numbers.
    Then $\finiteFunctionJoin{m}{n}{A}\in
        \bijections{\funs{m}{A}\times\funs{n}{A}}{\funs{m+n}{A}}$.
\end{corollary}
\begin{proof}
    Follows by \auto{finite_function_split_bijection,finite_function_split_converse,
        bijection_converse_is_bijection}.
\end{proof}

\begin{corollary}\label{finite_function_join_converse}
    Let $m,n$ be natural numbers.
    Then $\converse{\finiteFunctionJoin{m}{n}{A}}=\finiteFunctionSplit{m}{n}{A}$.
\end{corollary}
\begin{proof}
    $\finiteFunctionSplit{m}{n}{A}$ is a relation by
        \auto{finite_function_split_bijection,bijections_to_funs,funs_is_relation}.
    Follows by \auto{finite_function_split_converse,converse_converse_eq}.
\end{proof}

\begin{theorem}\label{finite_function_join_split}
    Let $m,n$ be natural numbers.
    Then $\finiteFunctionJoin{m}{n}{A}\circ\finiteFunctionSplit{m}{n}{A}
        =\identity{\funs{m+n}{A}}$.
\end{theorem}
\begin{proof}
    Follows by \auto{finite_function_split_bijection,finite_function_split_converse,
        bijection_converse_circ}.
\end{proof}

\begin{theorem}\label{finite_function_split_join}
    Let $m,n$ be natural numbers.
    Then $\finiteFunctionSplit{m}{n}{A}\circ\finiteFunctionJoin{m}{n}{A}
        =\identity{\funs{m}{A}\times\funs{n}{A}}$.
\end{theorem}
\begin{proof}
    Follows by \auto{finite_function_split_bijection,finite_function_split_converse,
        bijection_circ_converse}.
\end{proof}

\subsection{The coordinate formulas}

\begin{theorem}\label{finite_function_split_apply}
    Let $m,n$ be natural numbers and $x\in\funs{m+n}{A}$.
    Then $(\finiteFunctionSplit{m}{n}{A})(x)
        =(\restrl{x}{m},x\circ\finiteIndexShift{m}{n})$.
\end{theorem}
\begin{proof}
    Let $D=(m+n)\setminus m$ and $e=\finiteIndexShift{m}{n}$.
    Let $r=\functionReindex{D}{n}{A}{e}$ and $s=\functionSplit{m}{D}{A}$.
    Let $P=\productMap{\funs{m}{A}}{\funs{D}{A}}{\identity{\funs{m}{A}}}{r}$.
    $m\subseteq m+n$ by \obv{natural_addition_left_subseteq_sum}.
    $D\subseteq m+n$ by \auto{setminus_subseteq}.
    $m\union D=m+n$ by \auto{setminus_partition}.
    $e\in\funs{n}{D}$ by \auto{finite_index_shift_type}.
    Hence $e\in\funs{n}{m+n}$ by \obv{funs_weaken_codom}.
    $r\in\funs{\funs{D}{A}}{\funs{n}{A}}$ by \auto{function_reindex_type}.
    $\identity{\funs{m}{A}}\in\funs{\funs{m}{A}}{\funs{m}{A}}$ by \obv{id_elem_funs}.
    $P\in\funs{\funs{m}{A}\times\funs{D}{A}}{\funs{m}{A}\times\funs{n}{A}}$ by
        \auto{product_map_type}.
    $s\in\funs{\funs{m+n}{A}}{\funs{m}{A}\times\funs{D}{A}}$ by
        \auto{function_split_type}.
    $s(x)=(\restrl{x}{m},\restrl{x}{D})$ by \auto{function_split_apply}.
    $\restrl{x}{m}\in\funs{m}{A}$ and $\restrl{x}{D}\in\funs{D}{A}$ by \obv{funs_restrl}.
    $r(\restrl{x}{D})=(\restrl{x}{D})\circ e$ by \auto{function_reindex_apply}.
    $s(x)\in\funs{m}{A}\times\funs{D}{A}$ by \obv{funs_type_apply}.
    Hence $P(s(x))=(\restrl{x}{m},(\restrl{x}{D})\circ e)$ by
        \auto{product_map_apply,id_apply,fst_eq,snd_eq}.
    $(\restrl{x}{D})\circ e\in\funs{n}{A}$ and $x\circ e\in\funs{n}{A}$ by
        \obv{funs_circ}.
    For every $k\in n$ we have $((\restrl{x}{D})\circ e)(k)=(x\circ e)(k)$ by
        \auto{circ_apply,restrl_of_function_apply,funs_type_apply,funs_is_function,funs,funs_ran}.
    Hence $(\restrl{x}{D})\circ e=x\circ e$ by
        \obv{typed_functions_equal_from_agreement}.
    $(P\circ s)(x)=P(s(x))$ by \auto{circ_apply,funs_is_function,funs,funs_ran}.
    Follows by \auto{finite_function_split}.
\end{proof}

\begin{theorem}\label{finite_function_join_apply}
    Let $m,n$ be natural numbers and $p\in\funs{m}{A}\times\funs{n}{A}$.
    Then $(\finiteFunctionJoin{m}{n}{A})(p)
        =\fst{p}\union((\snd{p})\circ\converse{\finiteIndexShift{m}{n}})$.
\end{theorem}
\begin{proof}
    Let $D=(m+n)\setminus m$ and $e=\finiteIndexShift{m}{n}$.
    Let $r=\functionReindex{n}{D}{A}{\converse{e}}$ and $j=\functionJoin{m}{D}{A}$.
    Let $P=\productMap{\funs{m}{A}}{\funs{n}{A}}{\identity{\funs{m}{A}}}{r}$.
    $e\in\bijections{n}{D}$ by \auto{finite_index_shift_bijection}.
    Hence $\converse{e}\in\funs{D}{n}$ by
        \obv{bijection_converse_is_bijection,bijections_to_funs}.
    $r\in\funs{\funs{n}{A}}{\funs{D}{A}}$ by \auto{function_reindex_type}.
    $\identity{\funs{m}{A}}\in\funs{\funs{m}{A}}{\funs{m}{A}}$ by \obv{id_elem_funs}.
    $P\in\funs{\funs{m}{A}\times\funs{n}{A}}{\funs{m}{A}\times\funs{D}{A}}$ by
        \auto{product_map_type}.
    $\fst{p}\in\funs{m}{A}$ and $\snd{p}\in\funs{n}{A}$ by \obv{fst_type,snd_type}.
    $r(\snd{p})=(\snd{p})\circ\converse{e}$ by \auto{function_reindex_apply}.
    $P(p)=(\fst{p},(\snd{p})\circ\converse{e})$ by \auto{product_map_apply,id_apply}.
    $P(p)\in\funs{m}{A}\times\funs{D}{A}$ by \obv{funs_type_apply}.
    Hence $j(P(p))=\fst{p}\union((\snd{p})\circ\converse{e})$ by
        \auto{function_join_apply,fst_eq,snd_eq}.
    $m$ is disjoint from $D$ by \auto{disjoint,setminus}.
    $j\in\funs{\funs{m}{A}\times\funs{D}{A}}{\funs{m\union D}{A}}$ by
        \auto{function_join_type}.
    $(j\circ P)(p)=j(P(p))$ by \auto{circ_apply,funs_is_function,funs,funs_ran}.
    Follows by \auto{finite_function_join}.
\end{proof}

\begin{theorem}\label{finite_function_join_components}
    Let $m,n$ be natural numbers and $p\in\funs{m}{A}\times\funs{n}{A}$.
    Then $\restrl{(\finiteFunctionJoin{m}{n}{A})(p)}{m}=\fst{p}$ and
        $((\finiteFunctionJoin{m}{n}{A})(p))\circ\finiteIndexShift{m}{n}=\snd{p}$.
\end{theorem}
\begin{proof}
    Let $s=\finiteFunctionSplit{m}{n}{A}$ and $j=\finiteFunctionJoin{m}{n}{A}$.
    $s\in\funs{\funs{m+n}{A}}{\funs{m}{A}\times\funs{n}{A}}$ and
        $j\in\funs{\funs{m}{A}\times\funs{n}{A}}{\funs{m+n}{A}}$ by
        \auto{finite_function_split_bijection,finite_function_join_bijection,bijections_to_funs}.
    $s\circ j=\identity{\funs{m}{A}\times\funs{n}{A}}$ by
        \auto{finite_function_split_join}.
    Hence $s(j(p))=p$ by \auto{circ_apply,id_apply,funs_is_function,funs,funs_ran}.
    $j(p)\in\funs{m+n}{A}$ by \obv{funs_type_apply}.
    Hence $s(j(p))=(\restrl{j(p)}{m},j(p)\circ\finiteIndexShift{m}{n})$ by
        \auto{finite_function_split_apply}.
    Follows by \auto{fst_eq,snd_eq}.
\end{proof}

\begin{proposition}\label{finite_function_join_left_coordinate}
    Let $m,n$ be natural numbers and $p\in\funs{m}{A}\times\funs{n}{A}$.
    Let $i\in m$.
    Then $\apply{(\finiteFunctionJoin{m}{n}{A})(p)}{i}=(\fst{p})(i)$.
\end{proposition}
\begin{proof}
    Let $x=(\finiteFunctionJoin{m}{n}{A})(p)$.
    $x\in\funs{m+n}{A}$ by
        \auto{finite_function_join_bijection,bijections_to_funs,funs_type_apply}.
    $\restrl{x}{m}=\fst{p}$ by \auto{finite_function_join_components}.
    $m\subseteq m+n$ by \obv{natural_addition_left_subseteq_sum}.
    Hence $i\in m+n$ by \obv{subseteq}.
    Follows by \auto{restrl_of_function_apply,funs_is_function,funs}.
\end{proof}

\begin{proposition}\label{finite_function_join_right_coordinate}
    Let $m,n$ be natural numbers and $p\in\funs{m}{A}\times\funs{n}{A}$.
    Let $k\in n$.
    Then $\apply{(\finiteFunctionJoin{m}{n}{A})(p)}{m+k}=(\snd{p})(k)$.
\end{proposition}
\begin{proof}
    Let $x=(\finiteFunctionJoin{m}{n}{A})(p)$ and $e=\finiteIndexShift{m}{n}$.
    $x\in\funs{m+n}{A}$ by
        \auto{finite_function_join_bijection,bijections_to_funs,funs_type_apply}.
    $e\in\funs{n}{(m+n)\setminus m}$ by \auto{finite_index_shift_type}.
    Hence $e\in\funs{n}{m+n}$ by \auto{funs_weaken_codom,setminus_subseteq}.
    $x\circ e=\snd{p}$ by \auto{finite_function_join_components}.
    $e(k)=m+k$ by \auto{finite_index_shift_apply}.
    Follows by \auto{circ_apply,funs_is_function,funs,funs_ran}.
\end{proof}

\subsection{Empty blocks}

\begin{proposition}\label{finite_function_split_zero_left_apply}
    Let $n$ be a natural number and $x\in\funs{n}{A}$.
    Then $(\finiteFunctionSplit{\emptyset}{n}{A})(x)=(\emptyset,x)$.
\end{proposition}
\begin{proof}
    $\emptyset$ is a natural number by \obv{omega_empty}.
    $\emptyset+n=n$ by \obv{num_addition_zero_left}.
    $\finiteIndexShift{\emptyset}{n}=\identity{n}$ by \obv{finite_index_shift_zero}.
    $x\circ\identity{n}=x$ by \obv{funs_circ_identity_right}.
    $\emptyset\subseteq n$ by \obv{emptyset_subseteq}.
    $\restrl{x}{\emptyset}\in\funs{\emptyset}{A}$ by \obv{funs_restrl}.
    Hence $\restrl{x}{\emptyset}=\emptyset$ by
        \auto{empty_domain_function_space,singleton_iff}.
    Follows by \auto{finite_function_split_apply}.
\end{proof}

\begin{proposition}\label{finite_function_split_zero_right_apply}
    Let $n$ be a natural number and $x\in\funs{n}{A}$.
    Then $(\finiteFunctionSplit{n}{\emptyset}{A})(x)=(x,\emptyset)$.
\end{proposition}
\begin{proof}
    $\emptyset$ is a natural number by \obv{omega_empty}.
    $n+\emptyset=n$ by \obv{num_addition_zero_right}.
    $\finiteIndexShift{n}{\emptyset}\in\funs{\emptyset}{n\setminus n}$ by
        \auto{finite_index_shift_type}.
    Hence $\finiteIndexShift{n}{\emptyset}\in\funs{\emptyset}{n}$ by
        \auto{funs_weaken_codom,setminus_subseteq}.
    $x\circ\finiteIndexShift{n}{\emptyset}\in\funs{\emptyset}{A}$ by \obv{funs_circ}.
    Hence $x\circ\finiteIndexShift{n}{\emptyset}=\emptyset$ by
        \auto{empty_domain_function_space,singleton_iff}.
    $\restrl{x}{n}=x$ by \auto{restrl_by_dom,funs_is_relation,funs}.
    Follows by \auto{finite_function_split_apply}.
\end{proof}

\begin{proposition}\label{finite_function_join_zero_left_apply}
    Let $n$ be a natural number and $x\in\funs{n}{A}$.
    Then $(\finiteFunctionJoin{\emptyset}{n}{A})((\emptyset,x))=x$.
\end{proposition}
\begin{proof}
    Let $s=\finiteFunctionSplit{\emptyset}{n}{A}$ and $j=\finiteFunctionJoin{\emptyset}{n}{A}$.
    $\emptyset$ is a natural number by \obv{omega_empty}.
    $\emptyset+n=n$ by \obv{num_addition_zero_left}.
    $s\in\funs{\funs{n}{A}}{\funs{\emptyset}{A}\times\funs{n}{A}}$ and
        $j\in\funs{\funs{\emptyset}{A}\times\funs{n}{A}}{\funs{n}{A}}$ by
        \auto{finite_function_split_bijection,finite_function_join_bijection,bijections_to_funs}.
    $j\circ s=\identity{\funs{n}{A}}$ by \auto{finite_function_join_split}.
    Hence $j(s(x))=x$ by \auto{circ_apply,id_apply,funs_is_function,funs,funs_ran}.
    Follows by \auto{finite_function_split_zero_left_apply}.
\end{proof}

\begin{proposition}\label{finite_function_join_zero_right_apply}
    Let $n$ be a natural number and $x\in\funs{n}{A}$.
    Then $(\finiteFunctionJoin{n}{\emptyset}{A})((x,\emptyset))=x$.
\end{proposition}
\begin{proof}
    Let $s=\finiteFunctionSplit{n}{\emptyset}{A}$ and $j=\finiteFunctionJoin{n}{\emptyset}{A}$.
    $\emptyset$ is a natural number by \obv{omega_empty}.
    $n+\emptyset=n$ by \obv{num_addition_zero_right}.
    $s\in\funs{\funs{n}{A}}{\funs{n}{A}\times\funs{\emptyset}{A}}$ and
        $j\in\funs{\funs{n}{A}\times\funs{\emptyset}{A}}{\funs{n}{A}}$ by
        \auto{finite_function_split_bijection,finite_function_join_bijection,bijections_to_funs}.
    $j\circ s=\identity{\funs{n}{A}}$ by \auto{finite_function_join_split}.
    Hence $j(s(x))=x$ by \auto{circ_apply,id_apply,funs_is_function,funs,funs_ran}.
    Follows by \auto{finite_function_split_zero_right_apply}.
\end{proof}

\subsection{Three consecutive blocks}

\begin{theorem}\label{finite_function_join_associative}
    Let $m,n,k$ be natural numbers.
    Let $f\in\funs{m}{A}$ and $g\in\funs{n}{A}$.
    Let $h\in\funs{k}{A}$.
    Then $(\finiteFunctionJoin{m+n}{k}{A})(((\finiteFunctionJoin{m}{n}{A})((f,g)),h))
        =(\finiteFunctionJoin{m}{n+k}{A})((f,(\finiteFunctionJoin{n}{k}{A})((g,h))))$.
\end{theorem}
\begin{proof}
    Let $u=(\finiteFunctionJoin{m}{n}{A})((f,g))$ and
        $v=(\finiteFunctionJoin{n}{k}{A})((g,h))$.
    Let $l=(\finiteFunctionJoin{m+n}{k}{A})((u,h))$ and
        $r=(\finiteFunctionJoin{m}{n+k}{A})((f,v))$.
    $m+n$ and $n+k$ are natural numbers by \obv{num_addition_on_naturals}.
    $(m+n)+k=m+(n+k)$ by \obv{natural_addition_associative}.
    $(f,g)\in\funs{m}{A}\times\funs{n}{A}$ and
        $(g,h)\in\funs{n}{A}\times\funs{k}{A}$ by \obv{times_tuple_intro}.
    $u\in\funs{m+n}{A}$ and $v\in\funs{n+k}{A}$ by
        \auto{finite_function_join_bijection,bijections_to_funs,funs_type_apply}.
    $(u,h)\in\funs{m+n}{A}\times\funs{k}{A}$ and
        $(f,v)\in\funs{m}{A}\times\funs{n+k}{A}$ by \obv{times_tuple_intro}.
    $l\in\funs{m+(n+k)}{A}$ and $r\in\funs{m+(n+k)}{A}$ by
        \auto{finite_function_join_bijection,bijections_to_funs,funs_type_apply}.
    Show that for every $i\in m+(n+k)$ we have $l(i)=r(i)$.
    \begin{subproof}
        Fix $i\in m+(n+k)$.
        $i\in m$ or there exists $j\in n+k$ such that $i=m+j$ by
            \obv{natural_sum_member_iff}.
        \begin{byCase}
            \caseOf{$i\in m$.}
                $m\subseteq m+n$ by \obv{natural_addition_left_subseteq_sum}.
                Hence $i\in m+n$ by \obv{subseteq}.
                $l(i)=u(i)$ and $r(i)=f(i)$ by
                    \auto{finite_function_join_left_coordinate,fst_eq}.
                $u(i)=f(i)$ by \auto{finite_function_join_left_coordinate,fst_eq}.
            \caseOf{There exists $j\in n+k$ such that $i=m+j$.}
                Take $j\in n+k$ such that $i=m+j$.
                $r(i)=v(j)$ by \auto{finite_function_join_right_coordinate,snd_eq}.
                $j\in n$ or there exists $t\in k$ such that $j=n+t$ by
                    \obv{natural_sum_member_iff}.
                \begin{byCase}
                    \caseOf{$j\in n$.}
                        $i\in m+n$ by \auto{natural_sum_member_iff}.
                        $l(i)=u(i)$ by \auto{finite_function_join_left_coordinate,fst_eq}.
                        $u(i)=g(j)$ by \auto{finite_function_join_right_coordinate,snd_eq}.
                        $v(j)=g(j)$ by \auto{finite_function_join_left_coordinate,fst_eq}.
                    \caseOf{There exists $t\in k$ such that $j=n+t$.}
                        Take $t\in k$ such that $j=n+t$.
                        $k\subseteq\naturals$ by
                            \auto{naturals_is_transitive_set,transitiveset_iff_subseteq}.
                        Hence $t$ is a natural number by \obv{subseteq}.
                        $i=(m+n)+t$ by \auto{natural_addition_associative}.
                        $l(i)=h(t)$ and $v(j)=h(t)$ by
                            \auto{finite_function_join_right_coordinate,snd_eq}.
                \end{byCase}
        \end{byCase}
    \end{subproof}
    Follows by \auto{typed_functions_equal_from_agreement}.
\end{proof}
